% !TEX encoding = ISO-8859-1
\documentclass[12pt]{article}

\usepackage{amssymb,amsmath,amsfonts,amsthm,amstext}
\usepackage[brazil]{babel}
\usepackage[latin1]{inputenc}
\usepackage[pdftex]{graphicx}
\usepackage{enumitem}
\usepackage{multicol}
\usepackage[all]{xy}

\setlength{\topmargin}{-1.0in}
\setlength{\oddsidemargin}{0in}
\setlength{\textheight}{10.1in}
\setlength{\textwidth}{6.5in}
\setlength{\baselineskip}{12mm}

\newcounter{exercicios}
\setcounter{exercicios}{0}
\newcommand{\questao}{
\addtocounter{exercicios}{1}
\noindent{\bf Exerc{\'\i}cio \arabic{exercicios}: }}

\newcommand{\equi}{\Leftrightarrow}
\newcommand{\bic}{\leftrightarrow}
\newcommand{\cond}{\rightarrow}
\newcommand{\impl}{\Rightarrow}
\newcommand{\nao}{\sim}
\newcommand{\sub}{\subseteq}
\newcommand{\e}{\ \wedge\ }
\newcommand{\ou}{\ \vee\ }
\newcommand{\vaz}{\emptyset}
\renewcommand{\sin}{{\rm sen\,}}

\newcommand{\real}{\mathbb{R}}
\newcommand{\vesp}{\vspace{0.2cm}}
\newcommand{\z}{\mathbb{Z}}
\newcommand{\n}{\mathbb{N}}
\newcommand{\q}{\mathbb{Q}}
\newcommand{\subne}{\subsetneqq}


\newcommand{\compcent}[1]{\vcenter{\hbox{$#1\circ$}}}
\newcommand{\comp}{\mathbin{\mathchoice
{\compcent\scriptstyle}{\compcent\scriptstyle}
{\compcent\scriptscriptstyle}{\compcent\scriptscriptstyle}}}

\begin{document}
\pagestyle{empty}

\begin{figure}[h]
        \begin{minipage}[c]{1.7cm}
        \includegraphics[width=1.7cm]{../../../imagens/unb.pdf}
        \end{minipage}%
        \hspace{0pt}
        \begin{minipage}[c]{4in}
          {Universidade de Bras{\'\i}lia} \\
          {Departamento de Matem{\'a}tica}
\end{minipage}
\end{figure}
\vspace{-1cm}\hrule

\begin{center}
{\Large\bf {\'A}lgebra 1 - Turma A} \\ \vspace{9pt} {\large\bf
  $2^{\underline{a}}$ Lista de Exerc{\'\i}cios -- 2$^{o}$/2011}\\ \vspace{9pt} Prof. Jos{\'e} Ant{\^o}nio O. Freitas
\end{center}
\hrule

\vspace{.6cm}


\questao Quais das rela{\c c}{\~o}es abaixo s{\~a}o rela{\c c}{\~o}es de equival{\^e}ncia sobre $E = \{a,b,c\}$?
\begin{enumerate}[label={\alph*})]
\item $R_1 = \{(a,a),(a,b),(b,a),(b,b),(c,c)\}$;
\item $R_2 = \{(a,a),(a,b),(b,a),(b,b),(b,c)\}$;
\item $R_3 = \{(a,a),(b,b),(b,c),(c,b),(a,c),(c,a)\}$;
\item $R_4 = E \times E$;
\item $R_5 = \vaz$.
\end{enumerate}
\questao Determinar todas as rela{\c c}{\~o}es de equival{\^e}ncia
$R$ sobre $A$ e os respectivos conjuntos quocientes $A/R$ para:
\begin{enumerate}[label={\alph*})]
\item $A=\{a\}$;
\item $A=\{a,b\}$;
\item $A=\{a,b,c\}$;
\item $A=\{a,b,c,d\}$.
\end{enumerate}

\vesp

\questao Quais das seguintes senten{\c c}as definem uma rela{\c c}{\~a}o de equival{\^e}ncia em $\n$?
\begin{enumerate}[label={\alph*})]
\item $aRb$ se, e s{\'o} se, existe $k \in \n$ tal que $a - b = 3k$.
\item $a \mid b$.
\item $a \le b$.
\item $mdc(a, b) = 1$.
\item $x + y =10$.
\end{enumerate}
\vesp

\questao Seja $E$ um conjunto n{\~a}o vazio. Dados $X$, $Y \in \mathcal{P}(E)$, o conjunto das partes de $E$, mostre que as rela{\c c}{\~o}es $R$ e $S$ s{\~a}o de equival{\^e}ncia em  $\mathcal{P}(E)$.
\begin{enumerate}[label={\alph*})]
\item $X~R~Y \Leftrightarrow X\cap A=Y\cap A$.
\item $X~R~Y \Leftrightarrow X\cup A=Y\cup A$.
\end{enumerate}
onde $A$ {\'e} um subconjunto fixo de $E$.

\vesp

\questao Seja $A=\n\times \n^*$. Considere a seguinte
rela{\c c}{\~a}o sobre $A$:
\[
(a,b)R (c,d) \Leftrightarrow a + b = c + d.
\]
Mostre que $R$ {\'e} uma rela{\c c}{\~a}o de equival{\^e}ncia sobre $A$.

\vesp

\questao Seja $A=\z\times \z^*$, onde
$\mathbb{Z}^*=\mathbb{Z}\setminus \{0\}$. Para $(a,b), (c,d) \in
A$, considere a seguinte rela{\c c}{\~a}o
\[
(a,b)R (c,d) \Leftrightarrow ad=bc.
\]
\begin{enumerate}[label={\alph*})]
\item Mostre que $R$ {\'e} uma rela{\c c}{\~a}o de equival{\^e}ncia sobre $A$.
\item Descreva a classe de equival{\^e}ncia $\overline{(0,1)}$, $\overline{(1,1)}$, $\overline{(1,2)}$, $\overline{(2,1)}$, $\overline{(2,2)}$, $\overline{(2,3)}$.
\end{enumerate}

\vesp

\questao Considere a seguinte rela{\c c}{\~a}o sobre $\mathbb{C}$:
\[
(x+yi)R(r+si) \Leftrightarrow x^2+y^2=r^2+s^2.
\]
\begin{enumerate}[label={\alph*})]
\item Mostre que $R$ {\'e} rela{\c c}{\~a}o de equival{\^e}ncia.
\item Descreva a classe de equival{\^e}ncia de $1+i$.
\end{enumerate}

\vesp

\questao Seja $R$ uma rela{\c c}{\~a}o sobre $\q$
definida da seguinte forma:
\[
xRy \Leftrightarrow x-y \in \mathbb{Z}.
\]
\begin{enumerate}[label={\alph*})]
\item Prove que $R$ {\'e} uma rela{\c c}{\~a}o de equival{\^e}ncia sobre $\q$.
\item Descreva a classe $\bar{1}$.
\item Descreva a classe $\overline{1/2}$.
\end{enumerate}

\vesp

\questao A divisibilidade (ou seja, a rela{\c c}{\~a}o definida por $xRy$ se, e s{\'o}
se, $x \mid y$) {\'e} uma rela{\c c}{\~a}o de equival{\^e}ncia sobre $\z$?

\vesp

\questao Seja $R$ a seguinte rela{\c c}{\~a}o sobre $\z^*$:
\[
xRy \Leftrightarrow x\mid y \mbox{ e } y\mid x.
\]
Mostre que $R$ {\'e} uma rela{\c c}{\~a}o de equival{\^e}ncia sobre $\z^*$ e
descreva o conjunto quociente $\z^*/R$.

\vesp

\questao Seja $R = \{ (x, y) \in \real^2 \mid x - y \in \q\}$. Provar que $R$ {\'e} uma rela{\c c}{\~a}o de equival{\^e}ncia e descrever as classes representadas por $1/2$ e $\sqrt{2}$.

\vesp

\questao Seja $R$ uma rela{\c c}{\~a}o de equival{\^e}ncia
sobre um conjunto n{\~a}o vazio $E$. Considere o conjunto quociente
$E/R$ e  a aplica{\c c}{\~a}o quociente $q:E\to E/R$, ou seja, a aplica{\c c}{\~a}o
que a cada $x\in E$ associa sua classe $\bar{x}$ m{\'o}dulo $R$. Seja
$f:E\to F$ uma aplica{\c c}{\~a}o de $E$ em um conjunto $F$.

\begin{enumerate}[label={\alph*})]
\item Mostre que existe uma aplica{\c c}{\~a}o $g$ tal que $g\circ
q=f$ se, e somente se, a seguinte condi{\c c}{\~a}o estiver verificada:
para quaisquer $x,y\in E$, se $x\equiv y~(\mathrm{mod}~R)$ ent{\~a}o
$f(x)=f(y).$
\[
\xymatrix{
E  \ar[d]_-{f} \ar[r]^-{q}& E/R \ar@{-->}[dl]^-{g}  \\
F  &  \\
}
\]
\item Supondo que exista $g$ tal que $g\circ q=f$, mostre que
$g$ {\'e} {\'u}nica.
\item Nas condi{\c c}{\~o}es do item anterior, mostre que $g$ {\'e}
sobrejetora se, e somente se, $f$ {\'e} sobrejetora.
\end{enumerate}

\vesp

\questao Seja $f:E\to F$ e considere a seguinte
rela{\c c}{\~a}o sobre $E$:
\[
xRy \Leftrightarrow f(x)=f(y).
\]
\begin{enumerate}[label={\alph*})]
\item Mostre que $R$ {\'e} uma rela{\c c}{\~a}o de equival{\^e}ncia.
\item Mostre que se $f$ {\'e} sobrejetora, ent{\~a}o existe uma {\'u}nica
bije{\c c}{\~a}o $g: E/R \to F$ tal que $g\circ \pi = f$, onde $\pi: E\to
E/R$ {\'e} a aplica{\c c}{\~a}o quociente.
\end{enumerate}

\vesp

\questao Fa{\c c}a a tabela da soma e da multiplica{\c c}{\~a}o nos conjuntos: $\z_{3}$,
$\z_{4}$, $\z_{5}$, $\z_{6}$, $\z_{7}$, $\z_{8}$ e $\z_{9}$.

\vesp

\questao Resolva a equa{\c c}{\~a}o $x \oplus \overline{2} = \overline{0}$ no
conjunto $\z_{4}$. Tamb{\'e}m resolva as equa{\c c}{\~o}es $x \oplus \overline{2} =
\overline{1}$, $x^2 \oplus \overline{2}x = \overline{2}$ e $x^3 \oplus
\overline{2}x = \overline{3}$ no mesmo conjunto.

\vesp

\questao Para quais valores de $\overline{a}$ a equa{\c c}{\~a}o $\overline{2}\odot x = \overline{a}$ tem solu{\c c}{\~a}o no conjunto $\z_{5}$?

\vesp

\questao Resolva a equa{\c c}{\~a}o $x^2 \oplus \overline{2}x \oplus \overline{1} = \overline{0}$ nos conjuntos $\z_{8}$ e $\z_{9}$, caso ela tenha ra{\'\i}zes.

\vesp

\questao Mostre que para qualquer $n \in \z$, existe um elemento
$\overline{a} \in \z_{m}$, tal que $a \equiv n \pmod m$. ({\it Sugest{\~a}o:
  usar a divis{\~a}o de Euclides}.)

\vesp

\questao Para um n{\'u}mero inteiro positivo $m$ definimos a {\it fun{\c c}{\~a}o
  proje{\c c}{\~a}o} $f_m: \z \rightarrow \z_{m}$ por
\[
f_m(u) = \overline{u}
\]
onde $\overline{u}$ {\'e} a classe m{\'o}dulo $m$ de $u$ no conjunto
$\z_{m}$. Verifique se:
\begin{enumerate}[label=({\arabic*})]
\item $f_m$ {\'e} injetora?
\item $f_m$ {\'e} sobrejetora?
\item $f_m$ {\'e} bijetora?
\end{enumerate}

\vesp

\questao Com os dados do exerc{\'\i}cio anterior, mostre que
\[
f_m(x + y) = f_m(x) \oplus f_m(y),\quad f_m(xy) = f_m(x)\odot f_m(y).
\]

\vesp

\questao Por meio da rela{\c c}{\~a}o de congru{\^e}ncia m{\'o}dulo $m$, ache o resto da
divis{\~a}o do n{\'u}mero $7^{298}$ por $4$.

\vesp

\questao Por meio da rela{\c c}{\~a}o de congru{\^e}ncia m{\'o}dulo $m$, mostre que um n{\'u}mero
inteiro $n$ na base decimal {\'e} divis{\'\i}vel por $3$ se, e somente se, a soma de
seus algarismos {\'e} divis{\'\i}vel por $3$. Da mesma forma, mostre que um n{\'u}mero
positivo $n$ {\'e} divis{\'\i}vel por $5$ se, e somente se, o {\'u}ltimo algarismo de $n$
{\'e} $0$ ou $5$.

\vesp

\questao Mostre que nem sempre a soma de dois elementos de $\mathcal{U}(\z_{m})$ pertence ao conjunto  $\mathcal{U}(\z_m)$.

\vesp

\questao Ache todos os elementos invers{\'\i}veis de $\z_{100}$.

\vesp

\questao {\'E} falso ou verdadeiro que o n{\'u}mero de elementos dos conjuntos $\mathcal{U}(\z_{5})$ e $\mathcal{U}(\z_{10})^*$ {\'e} igual?

\vesp

\questao Mostre que para um n{\'u}mero primo $p$, h{\'a} somente duas possibilidades: $p \equiv 1 \pmod 4$ ou $p \equiv 3 \pmod 4$. {(\it Sugest{\~a}o: usar divis{\~a}o de Euclides por 4.)}

\vesp


\questao Consideremos em $\z \times \z$ as opera{\c c}{\~o}es $+$ e $\cdot$,
definidas por
\begin{enumerate}[label=({\arabic*})]
\item $(a, b) + (c, d) = (a + c, b + d)$,
\item $(a, b) \cdot (c, d) = (ac - bd, ad + bc)$.
\end{enumerate}
Mostre que $(\z\times \z, +, \cdot)$ {\'e} um anel comutativo e com unidade.

\vesp

\questao Prove que s{\~a}o an{\'e}is:
\begin{enumerate}[label=({\alph*})]
\item O conjunto $\z$ dotado da adi{\c c}{\~a}o usual e a multiplica{\c
    c}{\~a}o definida por $ab = 0$, para todos $a$, $b \in \z$.

\item O conjunto $\q$ com a soma $x \oplus y = x + y - 1$ e a multiplica{\c
    c}{\~a}o $x \odot y = x + y - xy$.

\item O conjunto $\z \times \z$ com a soma e multiplica{\c c}{\~a}o dadas por:
\[
(a,b) + (c,d) = (a + c, b + d) \quad \mbox{e} \quad (a,b) \cdot (c,d) = (ac, ad + bc).
\]
\end{enumerate}

\vesp

\questao Quais dos an{\'e}is do exerc{\'\i}cio anterior s{\~a}o comutativos?
Quais tem unidade? Determine a unidade no caso de existir.


\vesp

\questao Seja $M_2(\real)$ o conjunto das matrizes $2 \times 2$ com entradas
em $\real$. Defina a adi{\c c}{\~a}o e a multiplica{\c c}{\~a}o em
$M_2(\real)$ por
\[
\begin{pmatrix}
a & b\\
c & d
\end{pmatrix}
+
\begin{pmatrix}
e & f\\
g & h
\end{pmatrix}
=
\begin{pmatrix}
a + e & b + f\\
c + g & d + h
\end{pmatrix}
\]
\[
\begin{pmatrix}
a & b\\
c & d
\end{pmatrix}
\cdot
\begin{pmatrix}
e & f\\
g & h
\end{pmatrix}
=
\begin{pmatrix}
ae + bg & af + bh\\
ce + dg & cf + dh
\end{pmatrix}
\]
Prove que $(M_2(\real), +, \cdot)$ {\'e} um anel. Esse anel {\'e} comutativo? Possui unidade?

\vesp

\questao Considere a soma $+$ e a multiplica{\c c}{\~a}o $\cdot$ em $\mathbb{C}$ e $i$ o n�mero imagin�rio, isto �, $i^2 = -1$.
\begin{enumerate}[label=({\alph*})]
\item Seja $\z[i] = \{a + bi \mid a, b \in \z\}$. Mostre que $(\z[i], +,
  \cdot)$ {\'e} um anel comutativo com unidade. Esse anel {\'e} chamado {\it anel do
    inteiros de Gauss}.
\item Sejam $\z[\sqrt{3}] = \{ a + b\sqrt{3} \mid a, b \in \z\}$ e
  $\z[i\sqrt{3}] = \{ a + bi\sqrt{3} \mid a, b \in \z\}$. Mostre que
  $(\z[\sqrt{3}], +, \cdot)$ e $(\z[i\sqrt{3}], +, \cdot)$ s{\~a}o an{\'e}is comutativos
  com unidade.
\item Mais geralmente, mostre que se $n$ {\'e} um inteiro positivo, ent{\~a}o
  $(\z[\sqrt{n}], +, \cdot)$ e $(\z[i\sqrt{n}], +, \cdot)$ s{\~a}o an{\'e}is comutativos
  com unidade, onde $\z[i\sqrt{n}] = \{ a + bi\sqrt{n} \mid a, b \in \z\}$ e
  $\z[i\sqrt{n}] = \{ a + bi\sqrt{n} \mid a, b \in \z\}$ .
\end{enumerate}

\vesp

\questao Mostre que num anel comutativo $(A, +, \cdot)$ vale a identidade
$(a + b)^2 = a^2 + 2a\cdot b + b^2$. Como ser{\'a} escrita essa identidade nos
an{\'e}is como $\z_{2}$ e, em geral, nos an{\'e}is como $\z_{m}$, onde $m$ {\'e} um
n{\'u}mero par?

\vesp

\questao D{\^e} um exemplo de um anel com $4$ elementos todos satisfazendo a
condi{\c c}{\~a}o $x^2 = x$.

\vesp

\questao Dados um anel $(A, +, \cdot)$ e $n \in \n$. Para $a \in A$,
definimos
\[
a^1 = a \quad \mbox{e}\quad a^n = a^{n -1 }a
\]
para $n > 1$. Prove que
\begin{enumerate}[label=({\arabic*})]
\item Para todos $m$, $n \in \n$ e para todo $a \in A$, $a^ma^n = a^{m + n}$.
\item Para todos $m$, $n \in \n$ e para todo $a \in A$, $(a^m)^n = a^{mn}$.
\end{enumerate}

\vesp

\questao Seja $A$ um anel com a seguinte propriedade: $x^2 = x$, para todo
$x \in A$. Mostre que $-a = a$, para todo $a \in A$.

{\it Sugestao:
  Considere o produto $(x + y)^2$.}

\vesp

\questao Prove que um anel com a propriedade do exerc{\'\i}cio anterior �
necessariamente comutativo.

{\it Sugestao: Tome $x$, $y \in A$ e calcule $(x + y)^2$.}
\end{document}